\documentclass[10pt,a4paper]{amsart}
\usepackage[margin=19mm]{geometry}
\usepackage{amsmath,amssymb,mathtools}
\usepackage[scr=rsfs,cal=euler]{mathalfa}
\usepackage{xfrac}
\usepackage[unicode,hidelinks,bookmarks=false]{hyperref}
\newcommand{\ie}{i.e.\ }
\newcommand{\cf}{cf.\ }
\newcommand{\resp}{respectively\ }
\newcommand{\Subsection}{\subsection}
\providecommand{\nobreakdash}{\nobreak\hbox{-}}
\setlength{\emergencystretch}{2em}
\multlinegap0pt
\usepackage{amssymb}
\usepackage{xcolor}
\usepackage{tikz}
\newtheorem{lemma}{Lemma}
\usepackage{cleveref}
\crefname{lemma}{Lemma}{Lemmas}
\DeclareMathOperator{\col}{col}
\DeclareMathOperator{\row}{row}
\DeclareMathOperator{\upof}{up}
\title{Polynomial Time Enumeration of $t$-Stack-Sortable Permutations Ending in Their Least Entry}
\author{Jerry Zhang}
\date{Source: arXiv:2604.10779v3 (CC BY 4.0)}
\begin{document}
\maketitle

\noindent\emph{Source and licence.} Polynomial Time Enumeration of $t$-Stack-Sortable Permutations Ending in Their Least Entry. Version arXiv:2604.10779v3 (CC BY 4.0), distributed by arXiv under the Creative Commons Attribution 4.0 International licence, http://creativecommons.org/licenses/by/4.0/, as recorded in the arXiv licence metadata for this exact version and shown on the abstract page https://arxiv.org/abs/2604.10779v3. Full licence notice: \texttt{licenses/arxiv-2604.10779-CC-BY-4.0.txt}.\par\smallskip

\noindent\textit{Editorial scope.} Counted unit: the article's lemma lem:crux (source0.tex lines 699--703). The article's complete original proof of it is delivered with it (source0.tex lines 704--724, one \texttt{proof} environment of 21 lines). No mathematics is written, repaired or re-derived: the statement and the proof are the article's own lines, in the article's own notation, and no retained line is substituted or re-flowed.  Retained uncounted as the source-local context the proof needs: lines 674--677.  Nothing was omitted from any retained span, so the package carries no omission record.  No retained line contains a citation, so the wrapper carries no bibliography and no bibliography entry.  No retained line contains a figure environment, an image-inclusion command or drawing code, so no image is delivered and the package declares no asset dependency.  Editorial formatting only, in addition: an AMS \texttt{amsart} preamble that restates the article's own declarations and macros used by the retained lines, namely the environments \texttt{lemma} on separate counters, together with the standard packages those lines need.  The retained environments print in this document as Lemma 1 in order of appearance, whereas the article prints the same items with the numbering of its own full document.  The complete native source of the article as distributed in the arXiv e-print for this exact version is bundled unchanged as \texttt{sources/198292/source0.tex}.
\begin{lemma}\label{lem:b_alt_expression}
With the notation above, it holds that
\[B_{\pi,i,j_1,j_2}=\{T_{\pi}^{-1}(i',j')\mid (i',j')\in D(\alpha_\pi),i'\ge i, j_1<j'\le j_2\}.\]
\end{lemma}

\begin{lemma}\label{lem:crux}
With the notation above, the equality $T=T_{\pi}$ holds if and only if $(T^{-1}(1,j))_{j=1}^{\ell(\alpha)}$ is a subsequence of $\pi$ and
\[T^{-1}(i,j)\in H_{\pi,T}(i,j)\]
for all $(i,j)\in D(\alpha)$ such that $i>1$.
\end{lemma}

\begin{proof}
If $T=T_{\pi}$, then $(T^{-1}(1,j))_{j=1}^{\ell(\alpha)}$ is equal to $C_{\pi,1}$, which is a subsequence of $\pi$. Moreover, by \Cref{lem:b_alt_expression}, it holds that
\[T^{-1}(i,j) = T^{-1}_{\pi}(i,j)\in B_{\pi,1,\row_{\pi}(\upof_{\pi}(T_{\pi}^{-1}(i-1,j))),j}=H_{\pi,T}(i,j)\]
for all $(i,j)\in D(\alpha)$ such that $i>1$, which proves the forward direction.

Conversely, if $(T^{-1}(1,j))_{j=1}^{\ell(\alpha)}$ is a subsequence of $\pi$ and $T^{-1}(i,j)\in H_{\pi,T}(i,j)$ for all $(i,j)\in D(\alpha)$ such that $i>1$, we prove that $T^{-1}(i,j)=T^{-1}_{\pi}(i,j)$ by strong induction on $(i,j)$ in the poset $(D(\alpha),\ge_{D})$. For the maximal element $(i,j)=(1,1)$, we have
\[T^{-1}(1,1)=n=T_{\pi}^{-1}(1,1).\]
Now assume that the claim holds for all $(k,\ell)\in D(\alpha)$ such that either $(i,j)<_{D}(k,\ell)\le_{D}(1,1)$ or $k<i$. If $i=1$, then every element of $\{T^{-1}(i',j')\mid (i',j')\in D(\alpha), j'\le j-1\}$ is to the left of $T^{-1}(1,j-1)$ in $\pi$. Therefore, $T^{-1}(1,j)$ is the maximal element between $T^{-1}(1,j-1)=c_{\pi,1,j-1}$ and $0$ in $\pi$, which is precisely the definition of $c_{\pi,1,j}=T_{\pi}^{-1}(1,j)$. If $i>1$, then 
\[H_{\pi,T}(i,j)=B_{\pi,1,\row_{\pi}(\upof_{\pi}(T_{\pi}^{-1}(i-1,j))),j}\]
by the induction hypothesis because
\[T^{-1}(1,\row_{\alpha}(\upof_{\alpha}(i-1,j)))=T^{-1}_{\pi}(1,\row_{\pi}(\upof_{\pi}(T_{\pi}^{-1}(i-1,j))))\]
and
\[T^{-1}(1,j)=T^{-1}_{\pi}(1,j).\]
For all $1\le k\le i-1$, we also know that the last element of $B_{\pi,k,\row_{\pi}(\upof_{\pi}(T^{-1}_{\pi}(i-1,j))),j}$ is $T_{\pi}^{-1}(k,j)=T^{-1}(k,j)>T^{-1}(i,j)$, which implies that $\col_{\pi}(T^{-1}(i,j))$ must be at least $i$ and
\[T^{-1}(i,j)\in B_{\pi,i,\row_{\pi}(\upof_{\pi}(T^{-1}_{\pi}(i-1,j))),j}.\] Furthermore, no element of 
\[\{T^{-1}(i',j')\mid (i',j')\in D(\alpha), i'\le i-1\}\]
lies in $B_{\pi,i,\row_{\pi}(\upof_{\pi}(T^{-1}_{\pi}(i-1,j))),j}$ by \Cref{lem:b_alt_expression}, and similarly, no element of 
\[\{T^{-1}(i',j')\mid (i',j')\in D(\alpha), i'\ge i, j'<j\}\]
lies in $B_{\pi,i,\row_{\pi}(\upof_{\pi}(T^{-1}_{\pi}(i-1,j))),j}$ since any such element must be to the left of $T^{-1}_{\pi}(1,\row_{\pi}(\upof_{\pi}(T^{-1}_{\pi}(i-1,j))))$ in $\pi$. Therefore, we finally conclude that
\[T^{-1}(i,j)=\max(B_{\pi,i,\row_{\pi}(\upof_{\pi}(T^{-1}_{\pi}(i-1,j))),j})=T^{-1}_{\pi}(i,j).\]
\end{proof}

\end{document}
